\section{Analysis}
\label{sec:analysis}

This section separates observations stored in the analysis artifacts from their permitted interpretation.  It does not add a new system comparison, and it does not use error, confidence, or low-resource outputs to revise the main result ordering in Section~\ref{sec:results}.

\subsection{Confusion Analysis}

% TRACE: S6-C01 | answer/results/error_confusion.json | $.system, $.label, $.labels, $.row_normalised | exact matrix system and scope
\textbf{OBSERVED EVIDENCE.} The stored row-normalized pragmatic-polarity confusion matrix is for ViPragSent (ours), using its seed-20260520 prediction file, with six recorded classes: positive, negative, neutral, ironic-positive, sarcastic-negative, and mocking.  Rows are gold classes and columns are predicted classes.  It is therefore a single-seed diagnostic for this evaluated system and class construction, rather than an error analysis of the other systems or a three-seed aggregate from the main comparison.  The source graphic is reserved for an appendix placement so that its six class labels remain legible.
% TRACE: S6-C02 | answer/results/error_confusion.json | $.row_normalised[1][2], $.row_normalised[2][2], $.row_normalised[4][4], $.row_normalised[5][5] | descriptive matrix entries only
The recorded neutral diagonal is 79.76\%, whereas the negative row assigns 47.06\% of its records to neutral; the sarcastic-negative and mocking diagonals are 55.04\% and 55.08\%, respectively.
The ironic-positive row has no counts in this stored matrix, so it cannot support a row-level recall observation.  These values are descriptive entries in a row-normalized artifact, not independent estimates for every pragmatic label in the study.

\textbf{PLAUSIBLE INTERPRETATION.} The concentration of some negative and pragmatic-polarity records outside their diagonal cells identifies where this particular prediction set can be inspected more closely.  It supports a bounded statement that the recorded decision behavior is uneven across the represented classes, without establishing whether the source of an error is wording, annotation boundary, class prevalence, or the learned representation.

\textbf{UNTESTED HYPOTHESIS.} A targeted experiment could test whether rebalancing, alternative class construction, or a different prediction interface changes these off-diagonal patterns.  No such experiment is present in the artifact bundle, so these possibilities are not used as conclusions about the model architecture or the data.

\subsection{Calibration}

% TRACE: S6-C03 | answer/results/calibration.json | $.provenance.reporting_scope, $.systems.{phobert_finetune,xlmr_large,vipragsent_full} | eligible-system scope
\textbf{OBSERVED EVIDENCE.} Calibration is reported only for seed-20260520 prediction files with stored pragmatic-polarity confidence scores: PhoBERT fine-tune, XLM-R-large, and ViPragSent (ours).  The artifact excludes systems without those recorded scores, so it does not justify a calibration ranking across all seven main-comparison systems or a three-seed calibration estimate.
% TRACE: S6-C04 | answer/results/calibration.json | $.systems.{phobert_finetune,vipragsent_full,xlmr_large}.{n,ece,missing_confidence}; answer/tables/calibration.md | stored calibration values only
Each eligible system has 2,000 scored records and zero missing confidence entries; the stored ECE values are 0.09 for PhoBERT fine-tune, 0.07 for ViPragSent (ours), and 0.15 for XLM-R-large.
The bin-level reliability outputs and their visual detail are reserved for the appendix, where they can be read with the corresponding confidence-scope qualification.

\textbf{PLAUSIBLE INTERPRETATION.} Within this restricted set of recorded confidence outputs, ViPragSent has the lowest stored ECE.  This statement concerns calibration error on the defined pragmatic-polarity target; it neither changes the macro-F1 ordering nor implies that the same ordering would hold for systems without archived confidence values.

\textbf{UNTESTED HYPOTHESIS.} Whether calibration post-processing or a different confidence extraction procedure would alter these values is not tested here.  Accordingly, no causal account of the ECE differences is made.

\subsection{Exploratory Low-resource Sarcasm}

% TRACE: S6-C05 | answer/results/low_resource_sarcasm.json | $.provenance.reporting_scope, $.budgets.*.positive_budget, $.budgets.*.systems.*.metrics.sarcasm.n_seeds | exploratory one-seed scope
\textbf{OBSERVED EVIDENCE.} The low-resource sarcasm artifact is an exploratory comparison between PhoBERT fine-tune and ViPragSent (ours) at positive budgets of 64, 128, 256, 512, and 1024.  Each system--budget result has one registered seed, so the point values are not multi-seed estimates.
% TRACE: S6-C06 | answer/results/low_resource_sarcasm.json | $.budgets.{64,128,256,512,1024}.systems.*.metrics.sarcasm.mean | non-monotonic descriptive outcome only
The stored sarcasm F1 sequence for PhoBERT is 45.24, 44.96, 76.13, 75.99, and 76.33, while the ViPragSent sequence is 47.19, 48.72, 45.22, 74.31, and 74.49 across those budgets.
Both recorded sequences are non-monotonic: for example, the ViPragSent value falls from 48.72 at 128 to 45.22 at 256 before increasing at the later budgets.  The result is therefore reported as an exploratory pattern, not as evidence of data efficiency.

\textbf{PLAUSIBLE INTERPRETATION.} The stored curves show that budget alone does not provide a monotone description of these single-run outcomes.  They are useful for motivating a pre-registered multi-seed follow-up, but do not establish stable behavior across different sampled training subsets.

\textbf{UNTESTED HYPOTHESIS.} The observed changes could depend on the composition of the selected examples, optimization variation, or other unrecorded factors.  The present study does not isolate those alternatives, and it makes no architectural explanation or general low-resource claim.
